\documentclass[a4paper,12pt]{article}

\input{preamble}

\title{
    Linear Algebraic Analogues to the Graph Homomorphism Combinatorial Optimisation Problem
}
\author{
    Euan Mendoza
    \and
    Rachel Coster
}
% \date{June 1, 2023}

\begin{document}

\maketitle

\begin{abstract}
    Insert Abstract Here
\end{abstract}
\newpage

\tableofcontents

\section{Introduction}\label{sec:introduction}

The problem of deciding whether there exists a function preserving the edge relations of two graphs is known as the \textit{Graph Homomorphism Problem}. This problem in particular generalises some other more well known problems, such as the \textit{Graph Isomorphism Problem} ($\GIp$) and the \textit{Subgraph Isomorphism Problem}. The Subgraph Isomorphism Problem, also generalises well known problems such as the $k$-Clique Problem, and the Hamiltonian Path/Cycle problems.

The Graph Isomorphism Problem and the Subgraph Isomorphism Problems have been interesting problems that have been under investigation for as long as Cook and Levin's initial investigation into the class $\NP$. Graph Isomorphism is in the class $\NP \cap \coAM$ suggesting that Graph Isomorphism is not NP-complete unless the polynomial-time hierarchy collapses to the second level. On the other hand Subgraph Isomorphism (and by extension Graph Homomorphism) testing are canonical NP-complete problems, appearing in Karp's seminal paper\footnote{Karp also mentions Graph Isomorphism as an open question.}.

\todon{Practical complexity of graph isomorphism, Weisfeiler-Leman, Linear Algebraic Analogues to Weisfeiler-Leman. Also existing methods for subgraph isomorphism testing.}

\subsection{Preliminaries}\label{sec:preliminaries}

For $n\in \N$, let $[n] := \{ 1, 2, \ldots, n \}$.

\paragraph{Groups.} Given a graph $G$, $\Aut(G)$ denotes the group of automorphisms of a graph $G$. $\S_{n}$ denotes the \textit{symmetric group} of order $n$. $\pS(n)$ denotes the group of $n{\times}n$ \textit{permutation matrices}, matrices obtained by permuting the rows of the identity matrix. $\GL(n,\F)$ denotes the general linear group over an arbitrary field $\F$, clearly $\pS(n)$ is a subgroup of $\GL(n,\F)$.

\todon{write a description of }

\paragraph{Group Actions and equivalence relations.} Given matrices $A,B\in\M(n,\F)$

\section{Background to Graph Homomorphism}\label{sec:background}

\todom{Formal definitions of the problem and introduction of linear algebraic formulations.}

\section{The Algorithm}\label{sec:algorithm}

\section{Results}\label{sec:results}

\section{Conclusion}\label{sec:conclusion}

\end{document}
